\begingroup
\ArithmeticHead{242}{L}{ANALYSIS. BRIEFWECHSEL.}
\SourceLeaf{249b}
\pdfbookmark[1]{[4.] Schumacher an Gauss, 2. April 1808}{x1-correspondence-4}
\NumberedItem{4}
\begin{center}\textsc{Schumacher} an \textsc{Gauss}.\end{center}
\HistoricCitation{[Briefwechsel zwischen \textsc{Gauss} und \textsc{Schumacher} I, Altona 1860, S. 1.]}
\par\indent \{Sie werden sich, verehrter Herr Professor! über die Zuschrift eines Unbekannten wundern, vielleicht sie etwas zu dreist finden. Ehe Sie aber mich verurtheilen, bitte ich den Brief auszulesen; ich hoffe Sie werden dann selbst gestehen, dass ich in der Sache mich nur an Sie wenden konnte.

\par\indent Vor ungefähr 10 Jahren gab ein Spanier \textsc{Pedrayes} allen Mathematikern eine Differentialgleichung zu integriren auf, die in den Gött. Anz. und \textsc{Hindenburg}’s Archiv\textsuperscript{[*)]} abgedruckt ist. \textsc{Pfaff} hat sie nur für \textls[120]{einen besonderen} Fall integrirt (so viel ich noch erinnere setzt er die beyden Veränderlichen gleich), sonst ist meines Wissens gar nichts geschehen, so ungeheuer es auch scheinen mag, dass, da Sie, \textsc{Laplace}, \textsc{Lagrange} leben, ein Fremder solche Männer vergeblich ausfordern darf. Freilich sagt er, dass neue Methoden, Methoden die er erfunden, dazu erforderlich wären, aber kann er nicht durch Umwege erreicht haben, wozu man auf kürzerem Wege hätte kommen können? Und was er erfunden hat, kann es nicht nacherfunden werden? Ich gestehe Ihnen, wie ich voriges Jahr zufällig das Stück von \textsc{Hindenburg} in die Hand bekam, und daraus diese Umstände erfuhr, überlief mich ein Schauder, das Blut kochte mir, und in demselben Augenblick fasste ich den Entschluss (lächeln Sie immer) mich ganz der Mathematik zu widmen, zu der ich schon als Knabe mich hingerissen fühlte, und die ich seit der Zeit ohne Lehrer
\par\vfill\noindent\rule{.25\linewidth}{.4pt}\par\smallskip
\begingroup\footnotesize\noindent [*) Göttingische gelehrte Anzeigen 1798, 37. Stück, S. 361, \textsc{Hindenburgs} Archiv der r. u. a. Mathematik, Heft 9, 1799, S. 85.]\par\endgroup
\clearpage
\ArithmeticHead{243}{R}{SCHUMACHER AN GAUSS.}
\SourceLeaf{250}
\noindent durch Selbststudium verfolgt habe. — Ich hatte hier Gelegenheit mit dem Generale der spanischen Truppen, dem Marquis von \textsc{Romaña} bekannt zu werden, und habe von ihm das Versprechen erhalten, er werde mir \textsc{Pedrayes} eigne Auflösung schaffen.

\par\indent Kaum hatte ich das Versprechen erhalten, so dachte ich an Sie. Gewiss, Herr Professor, Sie können die Aufgabe lösen, wenn Sie wollen! Wie ehrenvoll wäre es für Deutschland, wenn Sie \textls[120]{unser Stolz}, ehe die Auflösung vielleicht aus Spanien kömmt, sie hier gäben!

\par\indent Wollten Sie mich durch ein paar Zeilen benachrichtigen, ob Sie von Ihrer kostbaren Zeit, hiezu einige Stunden abbrechen wollen, so würden Sie mich unendlich verbinden.
\par\indent Mit der unbegränztesten Hochachtung
\begin{flushright}
Ihr ergebenster\par\medskip
\textsc{Schumacher},
\end{flushright}
\noindent\begin{minipage}[t]{.40\linewidth}Altona, den 2. April 1808.\end{minipage}\hfill\begin{minipage}[t]{.58\linewidth}\raggedleft
(Dr. der Rechte, Palmaillenstrasse\\
im Hause der Conferenzräthin \textsc{Schumacher}.)\}
\end{minipage}
\clearpage
\endgroup
